\setlength{\emergencystretch}{4em}
% Source printed p. 210 / PDF p. 218.
\chapter{稳定曲线的模空间}\label{src:6}
% Explicit XeLaTeX page destinations keep every terminology-index page clickable.
\hypersetup{pageanchor=false}
\AddToHook{shipout/foreground}[chapter-six-index-pages]{\put(0,0){\special{pdf:dest (page.\thepage) [@thispage /Fit]}}}
\begin{epigraph}
关键在于, 我喜欢地图, 也就是``世界地图'', ``海洋图'', ``星空图''意义下的地图! 我想, 吸引我研究这组问题的一个重要因素, 是希望为代数簇世界的某些部分绘制地图. 在 Oscar Zariski 的课堂上, 代数簇仿佛是有形可触的东西, 那么为什么不能像麦哲伦那样出发探险, 揭示它们的地理呢?
\epigraphauthor{David Mumford~\cite[vii 页]{Mum04}}
\end{epigraph}
本章证明定理~\hyperref[thm:A]{A}: 对 $g\geq2$, $\overline{\mathcal M}_g$ 是维数为 $3g-3$ 的光滑, 固有且不可约的 Deligne--Mumford 叠, 并具有射影粗模空间 $\overline{\mathcal M}_g\to\overline M_g$. 按照 \S\bookref{1.7.2-subsection} 概述的构造射影模空间的六步策略, 我们依次进行如下步骤.
\begin{enumerate}[(1)]
\item (代数性) 将 $\overline{\mathcal M}_g$ 表为所有曲线的叠 $\mathcal M_g^{\mathrm{all}}$ 的子叠, 并证明 $\mathcal M_g^{\mathrm{all}}$ 是 $\mathbb Z$ 上局部有限型的代数叠 (定理~\bookref{6.4.6}).
\item (稳定性的开放性) $\overline{\mathcal M}_g\subseteq\mathcal M_g^{\mathrm{all}}$ 是开子叠 (命题~\bookref{6.3.24}).
\item (有界性) $\overline{\mathcal M}_g$ 有限型 (定理~\bookref{6.4.16}).
\item (稳定约化) $\overline{\mathcal M}_g$ 固有 (定理~\bookref{6.5.23}).
\item (模空间的存在性) 应用 Keel--Mori 定理 (\bookref{5.4.6}), 得到粗模空间 $\overline{\mathcal M}_g\to\overline M_g$, 其中 $\overline M_g$ 是固有代数空间.
\item (射影性) 用 Kollár 丰沛性判据 (\bookref{6.9.2}), 证明对 $k\gg0$, $\det\pi_*(\omega_{\mathcal U_g/\overline{\mathcal M}_g}^{\otimes k})$ 丰沛, 其中 $\pi:\mathcal U_g\to\overline{\mathcal M}_g$ 是万有族 (推论~\bookref{6.9.5}).
\end{enumerate}
在此过程中, 我们建立结点曲线 (\S\bookref{6.2}) 和稳定曲线 (\S\bookref{6.3} 与 \S\bookref{6.6}) 的基本性质. 为应用于模理论, 必须对曲线族证明其中大多数性质, 这是本章论述中技术难度最高的方面之一. 在 \S\bookref{6.7} 中证明 $\overline{\mathcal M}_g$ 的不可约性, 并在 \S\bookref{6.8} 中概述 $\overline M_g$ 的 GIT 构造.
